\scanpage{5}

\mspage{449}

qu'il puisse avoir appris quelque chose de moi. Pourtant, chose remarquable, mon nom est prononcé trois fois dans cet éloge (nullement funèbre pour le coup) de Deligne(``troisième médaille Fields de L'IHES''). Et même dans une périphrase il est fait allusion, avec le vague de rigueur qui doit entourer chaque apparition de ma modeste personne, au fait que j'aurais ``construit la théorie de cohomologi\unsure{e} en géométrie sur un corps quelconque'' - et sûrement encore ``avec la plus grande généralité naturelle'', ça sent la grothendieckerie à plein nez (*). La citation complète du contexte vaut la peine d'être donnée, c'est un petit chef-d'oeuvre du genre:

``Partant de là \handwritten{[}théorie de Hodge classique\handwritten{]} et d'analogies \unsure{l|!}-adiques suggérées par Grothendieck \handwritten{[}on se demande où Gr. a trouvé le temps d'apprendre des choses aussi sérieuses, tout en rédigeant ses vingt volumes de plus grandes généralités naturelles\handwritten{]}, il \handwritten{[}Deligne\handwritten{]} a dégagé la notion de structure de Hodge mixte et en a muni la cohomologie de toute variété algébrique complexe. En cohomologie l-adique, donc \handwritten{[}?\handwritten{]} pour des variétés sur un corps fini, il a prouvé les conjectures de Weil, d'une difficulté proverbiale. Ce résultat a paru d'autant plus surprenant \handwritten{[}!!\handwritten{]} que Grothendieck, après avoir construit la théorie de cohomologie en géométrie sur un corps quelconque \handwritten{[}on se demande bien ce qu'il est e\unsure{n}core allé chercher là\handwritten{]}, avait ramené la conjecture restante \handwritten{[}???\handwritten{]} à une série de conjectures qui sont aujourd'hui aussi inabordables qu'alors.''

En clair, bien loin d'avoir contribué en quoi que ce soit à prouver ce surprenant résultat d'une difficulté si proverbiale, ces grothendieckeries-là (au nom à faire fuir le généraliste-naturaliste le plus endurci) ont été tout juste bonnes à nous encombrer encore de \mbox{c o n j e c t u r e s} comme de juste (il n'en fait jamais d'autres!) et inabordables ce qui plus est (on s'en serait douté), tout autant aujourd'hui que quand il a eu l'idée saugrenue de les faire.

\noindent\rule{.22\linewidth}{.4pt}

(*) (18 mai) Dans l'Eloge Funèbre, il est question de la ``grande attention'' que je portais à la terminologie. Dans l'utilisation d'expressions saugrenues comme ``la plus grande généralité naturelle'' ou ``la théorie de cohomologie en géométri\unsure{e} sur un corps quelconque'', je perçois clairement l'intention de tourner en dérision cette attention.

Le soin extrême que j'accorde aux noms donnés aux choses découle naturell\unsure{e}ment du respect que j'ai pour ces choses, dont le nom est censé exprimer l'esse\unsure{n}ce, ou du moins quelque aspect essentiel. Par les echos qui me parviennent, j'a\unsure{i} été choqué plus d'une fois par l'affectation de dédain qui aujourd'hui semble d\unsure{e} mise vis-à-vis de cette attitude de respect, dédain qui s'exprime parfois par l'usage de noms abracadabrants pour des notions importantes. Voir aussi à ce sujet la note ``La Perversité'' (n°76).
% archive sheet number: 4
